% META: {"id": "1NSI-ARCHOS-AM-EXTRAIT", "chapitre": "1NSI-ARCHITECTURE-OS", "type_objet": "amenagee", "capacites": ["P-ARCH-01A", "P-ARCH-03B", "P-ARCH-03C"], "capacites_codes": ["C1", "C4", "C5"], "derive_de": ["1NSI-ARCHOS-EX-001", "1NSI-ARCHOS-EX-004", "1NSI-ARCHOS-EX-005"], "profil_amenagement": ["SEQUENCAGE", "ALLEGEMENT_ENONCE", "REPONSE_FERMEE", "AERATION", "AMORCE_FOURNIE", "RETRAIT_DOUBLE_TACHE"], "mode_creation": "adaptation", "sources_inspiration": [], "status": "generated"}
\section*{Version aménagée --- Extrait Architecture et systèmes}
\textit{Consignes séquencées, mise en page aérée, énoncés allégés.}

\bigskip

\subsection*{Exercice 1 --- Les constituants de la machine (C1)}

\textbf{Étape 1.} Relie chaque rôle à son constituant. Écris la lettre.

\bigskip

\begin{center}
\begin{tabular}{|l|c|}
\hline
\textbf{Rôle} & \textbf{Lettre} \\
\hline
Exécute les instructions & \dotfill \\
\hline
Conserve les données pendant l'exécution & \dotfill \\
\hline
Conserve les données quand la machine est éteinte & \dotfill \\
\hline
\end{tabular}
\end{center}

\medskip
\textbf{A.} processeur \qquad \textbf{B.} mémoire vive \qquad \textbf{C.} disque

\bigskip

\textbf{Étape 2.} Dans l'architecture de von Neumann, les instructions et les
données sont-elles rangées au même endroit ?

\medskip
\begin{center}
\fbox{Oui, dans la même mémoire} \qquad \fbox{Non, dans deux mémoires séparées}
\end{center}

\bigskip
\bigskip

\subsection*{Exercice 2 --- Agir en ligne de commande (C4)}

\textbf{Étape 1.} Écris la commande qui affiche le dossier courant.

\medskip
\dotfill

\bigskip

\textbf{Étape 2.} Complète la suite. Une commande par ligne.

\begin{console}
$ ______                  # ou suis-je ?
/home/eleve
$ ______                  # que contient ce dossier ?
notes.txt  photos
$ mkdir ______            # creer un dossier « projet »
$ ______ notes.txt projet/   # copier le fichier dedans
\end{console}

\bigskip

\textbf{Étape 3.} Après un \texttt{mkdir} réussi, le terminal n'affiche rien.
Est-ce normal ?

\medskip
\begin{center}
\fbox{Oui, le silence signale la réussite} \qquad \fbox{Non, c'est une erreur}
\end{center}

\bigskip
\bigskip

\subsection*{Exercice 3 --- Lire des droits (C5)}

Un fichier confidentiel porte les permissions \texttt{rw-r-----}.

\bigskip

\textbf{Étape 1.} Découpe en trois groupes de trois. Recopie-les.

\bigskip

\begin{center}
\begin{tabular}{|l|c|}
\hline
\textbf{Groupe} & \textbf{Trois caractères} \\
\hline
Propriétaire & \dotfill \\
\hline
Groupe & \dotfill \\
\hline
Autres & \dotfill \\
\hline
\end{tabular}
\end{center}

\bigskip

\textbf{Étape 2.} Additionne dans chaque groupe. Rappel : $r = 4$, $w = 2$, $x = 1$.

\bigskip

\begin{center}
\begin{tabular}{|l|c|}
\hline
\textbf{Groupe} & \textbf{Chiffre} \\
\hline
Propriétaire & $4 + 2 + 0 = \mathbf{6}$ \\
\hline
Groupe & \dotfill \\
\hline
Autres & \dotfill \\
\hline
\end{tabular}
\end{center}

\bigskip

\textbf{Étape 3.} Écris la commande complète :

\medskip
\texttt{chmod }\,\dotfill\, \texttt{confidentiel.txt}

\bigskip

\textbf{Étape 4.} Les autres utilisateurs peuvent-ils lire ce fichier ?

\medskip
\begin{center}
\fbox{Oui} \qquad \fbox{Non}
\end{center}
